\scanpage{29}
\deleted{Il lui correspond alors $\varphi^{-1}$ un élément
\deleted{unique}\inserted{canonique}
de $\Phi$ \unsure{il forme donc}}

sur $1\otimes A$ ou sur $A\otimes1$. L'utilisant
\deleted{pour} écrire ces \unsure{deux|deux iso|isom.} de modules sur
$A\otimes1$ \unsure{$\Phi\rightleftarrows A$|$\Phi\simeq A$}, et \deleted{utilisant}\inserted{composant avec} $\mathrm{tr}$ sur $A$
\deleted{pour écrire $A^*\simeq\mathrm{dual}_Q$ de $A\simeq A$}, on
trouve une forme sur $\Phi$, i.e. un élément noté
$\check\varphi$. Procédant à droite au lieu de prendre
à gauche on trouve pareil.

On \deleted{\unsure{veut|peut}} utilise une \uline{symétrie} sur $\Phi$ qui
échange les deux structures de $A$-modules.
L'application $\overset{\circ}{\Phi}{}^{\deleted{\unsure{*}}}\longrightarrow\overset{\circ}{\Phi}^{*}$
commute néc. à cette symétrie.

Pour tt élément $\varphi$ de $\overset{\circ}{\Phi}^s$ [fixe par
symétrie] on trouve une bijection
$A\longrightarrow\Phi$, d'où sur $A$ une \deleted{involution}
application $u\leadsto u^\varphi$ linéaire, telle que
$(u^\varphi)^\varphi=u$, [déterminée mod \deleted{conjugaison}
\inserted{\uuline{composition} \deleted{\unsure{avec}} aut. int.}
par un élément $a$ de $A$ tel que $a^\varphi=a$, si
on change $\varphi$ en $\psi$].
\margin{\tikz[baseline=(a.base)]{\node[inner sep=2pt] (a) {\shortstack[l]{$\unsure{Z|2}$\\\unsure{$\mathord{/}_{\circ}$|6|$\delta$} (mais non\,: conjugaison\\ par !)}};\draw (a.center) ellipse [x radius=1.6cm,y radius=.9cm];}}
Si on fait la $=$ construction pour $\check\varphi$, on trouve une
involution sur $A^\circ=A$ (ensemblistement)
qui est la \deleted{\unsure{sa}}$\mathord{/}^{\circ}$.
% slanted
\margin{$\mathord{/}^{\circ}$ On a ici $(uv)^\varphi=v^\varphi u^\varphi$ (indépendant
des deux particulier de $\psi$)}
La structure de \inserted{\unsure{départ}} bimodule de
$\Phi$ devient la structure $aub^\varphi$, et $\varphi\leadsto\varphi^{-1}$ devient
\inserted{$u\leadsto u^{-1}$} \ldots
% archive sheet number: 28
